\documentclass[11pt]{article}
\usepackage{rabbithole}

\phasenumber{15}
\phasetitle{Being a maintainer}
\phasesubtitle{The job is mostly not code. Governance, review, saying no, and not burning out.}

\hypersetup{pdftitle={Rabbit Holes, Phase 15: Being a maintainer},
            pdfauthor={Accelerate, Manipal University Jaipur},
            pdfsubject={git and GitHub handbook}}

\begin{document}
\maketitlepage

\section{What the job actually is}

In Phase 3 the team leader merged pull requests. Every time they merged one,
everyone else's branch went stale and somebody had to rebase. That small annoyance
is the whole job in miniature: a maintainer's decisions are mostly decisions about
other people's time.

The surprise, for anyone who takes on a project people actually use, is how little
of the work is programming. The realistic split is closer to this:

\begin{description}
  \item[Deciding] What belongs in this project and what does not. The most valuable
    and least visible thing you do.
  \item[Reviewing] Reading other people's code carefully enough to be useful, and
    quickly enough to be worth submitting to.
  \item[Explaining] The same five questions, forever, to people who each ask once.
  \item[Triaging] Sorting the incoming pile into real, duplicate, unclear, and not
    our problem.
  \item[Releasing] Versioning, notes, and not breaking people.
  \item[Refusing] Politely, clearly, and often.
  \item[Coding] Whatever is left.
\end{description}

Nobody warns you about this, and it is the main reason maintainers quit.

\begin{why}
This phase exists because Accelerate runs an organisation with repositories that
outlive their authors' membership. The people who wrote \co{campus-mapper} will
graduate. Somebody will inherit it. Everything here is the difference between
inheriting a project and inheriting a mess.
\end{why}

\section{The files that answer questions for you}

Every question you answer twice should be answered in a file instead. GitHub has
specific names it recognises and surfaces automatically.

\vspace{0.5em}
\begin{tabularx}{\linewidth}{@{}L{4.2cm} X@{}}
\toprule
\rhtablehead{File} & \rhtablehead{Answers} \\
\midrule
\co{README.md} & What is this, why would I use it, how do I start \\
\co{CONTRIBUTING.md} & How do I submit a change and what will you reject \\
\co{CODE\_OF\_CONDUCT.md} & How people are expected to behave, and enforcement \\
\co{SECURITY.md} & How to report a vulnerability privately \\
\co{SUPPORT.md} & Where to ask questions that are not bugs \\
\co{LICENSE} & What others may legally do \\
\co{CODEOWNERS} & Who must review which files \\
\co{FUNDING.yml} & Adds a Sponsor button \\
\co{CITATION.cff} & How to cite this in a paper \\
\bottomrule
\end{tabularx}
\vspace{0.5em}

Put them in the repository root or in \co{.github/}. An organisation-level
\co{.github} repository provides defaults for every repository that lacks its own,
which is exactly why Accelerate has one.

\begin{note}
The highest leverage sentence in \co{CONTRIBUTING.md} is the one describing what you
will \emph{not} accept. Scope creep arrives as a well-written, working pull request
from someone who spent their weekend on it, and rejecting that is miserable for
everyone. A sentence up front prevents the whole situation.
\end{note}

\section{Licensing}

\begin{trap}
A public repository with no LICENSE file is \textbf{not} open source. Under
copyright law, code with no licence is all rights reserved. Nobody may legally use,
modify, or distribute it.

GitHub's terms let others view and fork within GitHub, and that is all. Every serious
company's legal team blocks dependencies with no licence. If you want people to use
your work, add a licence on day one.
\end{trap}

\vspace{0.5em}
\begin{tabularx}{\linewidth}{@{}L{2.5cm} X@{}}
\toprule
\rhtablehead{Licence} & \rhtablehead{In short} \\
\midrule
MIT & Do what you like, keep the notice, no warranty. Short and universally
understood. The default for student and club projects. \\
Apache 2.0 & Like MIT, plus an explicit patent grant and a requirement to state
changes. Preferred by companies for that patent clause. \\
BSD 3-Clause & Like MIT, plus you may not use the authors' names to endorse. \\
GPL v3 & Anyone distributing a derivative must release their source under GPL too.
Copyleft. \\
AGPL v3 & GPL, and running it as a network service counts as distribution. Closes
the cloud loophole. \\
MPL 2.0 & Copyleft per file. Modified files stay open, the larger work need not be. \\
CC BY-SA & For documentation and creative work, not code. \\
\bottomrule
\end{tabularx}
\vspace{0.5em}

\begin{shell}
gh repo create project --public --license mit
gh api repos/:owner/:repo --jq '.license.spdx_id'
\end{shell}

\subsection{Contributions and provenance}

When someone contributes, who owns it? By default they retain copyright and licence
it to you under the project's licence. Two mechanisms formalise it:

\textbf{DCO}, the Developer Certificate of Origin, is a statement that you wrote the
code and may submit it, asserted by signing off commits:

\begin{shell}
git commit -s -m "Add parser"
\end{shell}

That appends \co{Signed-off-by: Name <email>}. Lightweight, and enough for most
projects, including the Linux kernel.

\textbf{CLA}, a Contributor Licence Agreement, is a real contract, usually assigning
broader rights to a company. Heavier, and a genuine barrier: a first-time contributor
fixing a typo will not sign a legal document, and you lose that contribution.

For a student club, DCO or nothing. A CLA is a corporate tool for corporate reasons.

\section{Protecting the branch}

\begin{shell}
gh api -X PUT repos/:owner/:repo/branches/main/protection \
  -F required_pull_request_reviews[required_approving_review_count]=1 \
  -F required_pull_request_reviews[dismiss_stale_reviews]=true \
  -F required_status_checks[strict]=true \
  -F required_status_checks[contexts][]=test \
  -F enforce_admins=false \
  -F restrictions=null
\end{shell}

The options worth understanding:

\begin{description}
  \item[Required reviews] How many approvals before merging. One is usually right.
    Two on a project where a mistake is expensive.
  \item[Dismiss stale reviews] New commits invalidate old approvals. Turn this on.
    Without it, someone can get approval and then push something else.
  \item[Require branches up to date] The \co{strict} flag. Prevents a green pull
    request merging into a \co{main} that has moved. Correct, and it causes the
    constant rebasing that merge queues solve.
  \item[Required status checks] Which CI jobs must pass.
  \item[Enforce admins] Whether the rules apply to you. Leaving this off is a
    pragmatic escape hatch and a slow corrosion of the rules. Turn it on once the
    project is stable.
  \item[Require signed commits] Phase 10.
  \item[Require linear history] Blocks merge commits, forcing squash or rebase.
  \item[Require conversation resolution] All review comments resolved before merge.
\end{description}

\begin{note}
\textbf{Rulesets} are the newer system and are better. They can apply to several
repositories at once, target branches by pattern, be evaluated in a dry-run mode
before enforcement, and layer rather than override. If you are setting this up fresh
in an organisation, use rulesets rather than per-repository branch protection.
\end{note}

\subsection{CODEOWNERS}

\begin{terminal}
*                       @accelerate-muj/core-committee
/handbook/              @TheQMLGuy
/.github/workflows/     @TheQMLGuy @shashwat-deep
*.tex                   @TheQMLGuy
/phase-2/poems/
\end{terminal}

Last matching rule wins, unlike \co{.gitignore}. The empty last line means poem
submissions need no specific owner, which is what you want during a workshop with
thirty simultaneous pull requests.

Combined with ``require review from code owners'' in branch protection, the right
person is requested automatically and nothing merges without them.

\section{Reviewing well}

Code review is a skill and most people are never taught it. Reviewing badly is worse
than not reviewing.

\subsection{What to actually look for, in order}

\begin{enumerate}
  \item \textbf{Is this the right change at all?} Wrong-problem is the most expensive
    thing to catch late and the cheapest to catch here. Ask before reading closely.
  \item \textbf{Does it work?} Check it out and run it. \co{gh pr checkout 42}.
  \item \textbf{What happens at the edges?} Empty input, no network, concurrent runs,
    huge input.
  \item \textbf{Can the next person understand it?} Including the author in a year.
  \item \textbf{Is it tested?} Enough that a future change breaking it is caught.
  \item \textbf{Style.} Last, and ideally not by you. A formatter should own this.
\end{enumerate}

\subsection{How to say it}

\begin{trap}
The two failure modes are the rubber stamp and the pile-on. Both are common and both
drive contributors away.

Things that help, concretely:

\begin{itemize}
  \item \textbf{Say what is good.} Not flattery. Reviews that are exclusively
    negative read as hostile even when every point is correct.
  \item \textbf{Separate must from nice.} Prefix optional remarks with
    \co{nit:} or \co{optional:}. Without that signal every comment reads as blocking
    and a two-line change takes a week.
  \item \textbf{Ask rather than assert.} ``What happens if this is empty?'' beats
    ``this crashes on empty input'', because sometimes it does not and you were
    wrong.
  \item \textbf{Explain the why.} ``Use a set here'' teaches nothing. ``A set makes
    this lookup constant time, and this runs per request'' teaches once and forever.
  \item \textbf{Use suggested changes} for anything trivial. A one-click fix is
    kinder than a paragraph.
  \item \textbf{Review promptly or say you cannot.} A pull request sitting untouched
    for two weeks is the single most common reason people stop contributing. ``I
    cannot look at this until Friday'' costs you ten seconds and keeps them.
\end{itemize}
\end{trap}

\begin{shell}
gh pr diff 42
gh pr checkout 42
gh pr review 42 --comment --body "Two questions inline, nothing blocking"
gh pr review 42 --approve --body "Nice catch on the empty case"
gh pr review 42 --request-changes --body "See the note about the token"
\end{shell}

\subsection{Receiving review}

The other half, and it is a skill too. The reviewer is not attacking you. If a
reviewer misread your code, that is evidence the code is easy to misread, and the
useful response is usually a clearer name or a comment rather than an explanation in
the thread that nobody will ever read again.

Push back when you disagree. A review is a conversation, not a verdict. Just do it
about the technical point.

\section{Triage}

Incoming issues are not a to-do list. They are unsorted input, and about half are
not work.

A workable label scheme, kept small:

\begin{terminal}
type:    bug  feature  docs  question
status:  needs-info  confirmed  blocked  wontfix  duplicate
effort:  good-first-issue  help-wanted
\end{terminal}

\begin{shell}
gh issue list --label "status:needs-info" --json number,updatedAt
gh issue edit 42 --add-label "status:confirmed" --milestone "v1.1"
gh issue close 17 --reason "not planned" --comment "Duplicate of #12"
\end{shell}

\subsection{The triage question}

For every new issue, in order: \emph{Can I reproduce it? Is it already reported? Is
it in scope? Is it actionable as written?}

Anything failing the last test gets one request for specifics and a
\co{status:needs-info} label. If it goes quiet for a few weeks, close it kindly.
Closing is not rejection, and reopening is free.

\begin{trap}
Stale bots that close inactive issues automatically are widely used and widely
resented. From the maintainer's side it keeps the list honest. From the reporter's
side, being told by a robot that a real, unfixed bug is now closed for inactivity is
insulting, and it is a recurring source of public complaints about major projects.

If you use one, exclude confirmed bugs, give a long window, and write the message
like a person. Better: do not close on inactivity at all, and use a label plus a
saved search instead. An issue nobody has fixed is still true.
\end{trap}

\section{Saying no}

The hardest part of the job, and the one that decides whether your project stays
maintainable.

Every accepted feature is permanent. You will maintain it, document it, keep it
working through refactors, and break someone if you remove it. A feature used by four
percent of users costs you forever.

\subsection{Wording that works}

For an out-of-scope request:

\begin{quote}
\emph{Thanks for writing this up. I do not think this belongs in the core project,
because it only applies when you are using X and it would add a dependency for
everyone else. It would make a good extension, though, and I would happily link it
from the README.}
\end{quote}

For a pull request nobody asked for:

\begin{quote}
\emph{I appreciate the work here and I am not going to merge it, so let me explain
why rather than leaving it open. This changes the default behaviour, and there are
people relying on the current default. If you want to pursue it, open an issue and
we can talk about a migration path.}
\end{quote}

The pattern in both: decide, say so plainly, give the actual reason, offer a path
forward if there is one. Do not soften it into ambiguity.

\begin{trap}
The worst thing you can do is leave it open. An unanswered pull request is a person
waiting, checking, and slowly concluding they were ignored. A clear no on the day it
arrives is far kinder than silence for six months, and people generally respect it.

If you are not sure, say that: ``I need to think about this, give me a week.'' Then
actually do.
\end{trap}

\section{Releases and not breaking people}

Phase 10 covered the mechanics. The judgement is here.

\textbf{Compatibility is a promise.} Once someone depends on your project, changing
behaviour breaks their build at a time you chose and they did not. Semantic versioning
is how you communicate that in advance.

\textbf{Deprecate before removing.} The sequence that works:

\begin{enumerate}
  \item Release a version where the old way still works but warns, and the new way
    exists.
  \item Document the migration, with a before and after example.
  \item Leave it for at least one minor release, longer if the project is widely used.
  \item Remove it in the next MAJOR, and say so prominently in the notes.
\end{enumerate}

\textbf{Write release notes for humans.} A list of commit messages is not release
notes. Lead with what changed for the user, then what to do about it, and put the
full log at the bottom for anyone who wants it.

\section{Security}

\co{SECURITY.md} states which versions you support and how to report privately.

\begin{trap}
Never ask people to report vulnerabilities as public issues. The whole point is a
window to fix it before it is exploited.

GitHub has \textbf{private vulnerability reporting}, free, which gives you a private
thread and a place to draft an advisory and request a CVE:

\begin{shell}
gh api -X PUT repos/:owner/:repo/private-vulnerability-reporting
\end{shell}

Turn it on for anything anyone else runs.
\end{trap}

The process when a report arrives: acknowledge within a day or two even if only to
say you have seen it; agree a disclosure date with the reporter, ninety days being the
industry norm; fix privately using a draft security advisory, which gives you a
private fork; release the fix and publish the advisory at the same moment; credit the
reporter, because it costs nothing and it is the entire compensation most researchers
get.

\subsection{Dependencies}

\begin{terminal}
version: 2
updates:
  - package-ecosystem: "pip"
    directory: "/"
    schedule: {interval: "weekly"}
    open-pull-requests-limit: 5
    groups:
      dev-dependencies:
        patterns: ["pytest*", "ruff"]
\end{terminal}

Group updates, or Dependabot will file eleven pull requests a week and you will start
ignoring them, which is worse than not having it.

\section{Growing contributors}

A project with one maintainer has a bus factor of one. Growing that is a deliberate
activity.

\begin{description}
  \item[Label real starter issues] \co{good first issue} on something genuinely
    small, with the file named and the approach sketched. A vague hard task with a
    friendly label is worse than nothing.
  \item[Answer the first pull request properly] Someone's first contribution is
    disproportionately important. Review it that week. Be generous about small
    things.
  \item[Promote people] Someone who has landed several good changes and reviews
    others should be given commit access before they ask. Most people will never ask.
  \item[Write down the tacit knowledge] The reason for the odd workaround, the thing
    that always breaks in CI. \co{CONTRIBUTING.md}, or an \co{ARCHITECTURE.md}.
  \item[Let go of things] The instinct to fix it yourself because it is faster is
    correct in the short run and fatal in the long run.
\end{description}

\section{Governance}

Small projects run on ``whoever started it decides''. That works until it does not.

\begin{description}
  \item[BDFL] One person has final say. Efficient, clear, and a single point of
    failure in every sense.
  \item[Meritocracy or committer model] Contributors earn commit access over time,
    decisions by rough consensus. Most mid-sized projects.
  \item[Elected steering committee] Formal, elections, defined terms. Overhead worth
    paying once enough people depend on it. Accelerate's constitution is this.
  \item[Foundation] Neutral legal owner, for projects where several companies need
    assurance nobody can capture it.
\end{description}

Whatever you use, write it down. Unwritten governance is the founder's preferences,
and it is invisible until the first real disagreement, at which point everyone
discovers they believed something different.

\section{Burnout}

The predictable ending. Open source maintainer burnout is common enough to be
studied, and the pattern is consistent: growing obligation, static reward, no
boundary between the project and the rest of life.

What actually helps:

\begin{description}
  \item[Set office hours] Fixed times you look at the project. Outside them, do not.
    The alternative is a background hum of obligation that never stops.
  \item[Say the project is finished] Software can be done. ``Complete, accepting bug
    fixes only'' in the README is a legitimate and underused state.
  \item[Reduce the surface] Fewer features, fewer platforms, fewer promises. Every
    reduction is permanent relief.
  \item[Recruit before you need to] Phase 15's whole previous section.
  \item[Turn off notifications] Genuinely. Phase 12.
  \item[Remember what you owe] You published something for free. You do not owe
    anyone support, a fix, a response, or your weekend. Entitled demands are the
    other person being wrong.
\end{description}

\begin{trap}
The specific thing to watch for is a project that has become an obligation you
resent and cannot put down. The exits are all fine: hand it over, add maintainers,
mark it unmaintained, or archive it.

Archiving is honest and much better than abandoning it silently while people keep
filing issues.

\begin{shell}
gh repo archive old-project
\end{shell}

Say why in the README, and point at a fork if a good one exists.
\end{trap}

\section{Endings}

\textbf{Handing over.} Add them as an admin, let them merge things for a while, then
transfer. \co{Settings, Transfer ownership} moves it with issues, pull requests and
stars intact, and leaves a redirect. Never delete and re-upload, which destroys all
of that.

\textbf{Being forked.} Anyone may fork your project. If someone forks because they
disagree with your direction, that is the licence working as intended, and it is not
a betrayal. If the fork is better maintained than yours, the honest move is to link
it.

\textbf{Succession for a club.} Accelerate's specific case. Every repository should
name a current owner, and that name should be updated when officers change. A repo
whose owner graduated two years ago is unmaintained regardless of what the settings
say.

\checkyourself{
\begin{enumerate}
  \item What can people legally do with a public repository that has no LICENSE?
  \item Why turn on ``dismiss stale reviews''?
  \item Name three things that make a code review useful rather than discouraging.
  \item Why is leaving a pull request open worse than rejecting it?
  \item What are the four steps to remove a feature without breaking people?
  \item Someone reports a vulnerability in a public issue. What now?
  \item Your project has one maintainer and it is you. What is the actual risk, and
    what reduces it?
\end{enumerate}}

\vspace{1em}
{\color{rhborder}\rule{\linewidth}{0.8pt}}

\textbf{Next:} Phase 16, the last one, on the parts of git that only show up at
scale, and on where git came from, which explains several things that look strange
until you know.

\end{document}
